%------------------------------------------------------------------------------%
\begin{question}
  Encontre a inversa de
  \quad
  \(
    A =
    \begin{bmatrix}
      2 & 1 \\
      1 & 1
    \end{bmatrix}
  \)
\end{question}

%------------------------------------------------------------------------------%
\begin{resolution}{Matriz Aumentada}
  Matriz aumentada \([\,A\,|\,I\,]\)
  \pause
  \[
    \begin{amatrix}[2]{2}
      \emph{2} & \emph{1} & 1 & 0 \\
      \emph{1} & \emph{1} & 0 & 1
    \end{amatrix}
  \]

  \pause
  Queremos transformar a parte da esquerda na identidade $I_2$
\end{resolution}

%------------------------------------------------------------------------------%
\begin{resolution}{Escalonamento}
\begin{columns}
\column{0.3\linewidth}
  \[\;\;\;
    \begin{amatrix}[2]{2}
      2 & 1 & 1 & 0 \\
      1 & 1 & 0 & 1
    \end{amatrix}
  \]
  \pause
\column{0.7\linewidth}
  Trocamos as linhas
  \\ \pause
  \(
    L_1\leftrightarrow L_2
  \)
\end{columns}

\pause
\[
  \begin{amatrix}[2]{2}
    1 & 1 & 0 & 1 \\
    2 & 1 & 1 & 0
  \end{amatrix}
\]
\end{resolution}

%------------------------------------------------------------------------------%
\begin{resolution}{Escalonamento}
\begin{columns}
\column{0.3\linewidth}
  \[\;\;\;
    \begin{amatrix}[2]{2}
      1 & 1 & 0 & 1 \\
      2 & 1 & 1 & 0
    \end{amatrix}
  \]
  \pause
\column{0.7\linewidth}
  Eliminamos os elementos abaixo do primeiro pivô
  \\ \pause
  \(
    L_2 \leftarrow L_2 - 2 L_1
  \)
\end{columns}
\vspace{-\baselineskip}
\[
  \pause
  L'_2
  \;=\;
  L_2 - 2 L_1
  \pause
  \;=\;
  \begin{amatrix}[2]{2}
    2 & 1 & 1 & 0
  \end{amatrix}
  - 2
  \begin{amatrix}[2]{2}
    1 & 1 & 0 & 1
  \end{amatrix}
  \pause
  \;=\;
  \begin{amatrix}[2]{2}
    0 & -1 & 1 & -2
  \end{amatrix}
\]
\pause
\[
  \begin{amatrix}[2]{2}
    1 & 1  & 0 & 1  \\
    0 & -1 & 1 & -2
  \end{amatrix}
\]
\end{resolution}

%------------------------------------------------------------------------------%
\begin{resolution}{Escalonamento}
\begin{columns}
\column{0.3\linewidth}
  \[\;\;\;
    \begin{amatrix}[2]{2}
      1 & 1  & 0 & 1  \\
      0 & -1 & 1 & -2
    \end{amatrix}
  \]
  \pause
\column{0.7\linewidth}
  Tornando o segundo pivô igual a 1
  \\ \pause
  \(
    L_2 \leftarrow - L_2
  \)
\end{columns}
\vspace{-\baselineskip}
\[
  \pause
  L'_2
  \;=\;
   - L_2
  \pause
  \;=\;
  -
  \begin{amatrix}[2]{2}
    0 & -1 & 1 & -2
  \end{amatrix}
  \pause
  \;=\;
  \begin{amatrix}[2]{2}
    0 & 1 & -1 & 2
  \end{amatrix}
\]
\pause
\[
  \begin{amatrix}[2]{2}
    1 & 1 & 0  & 1 \\
    0 & 1 & -1 & 2
  \end{amatrix}
\]
\end{resolution}

%------------------------------------------------------------------------------%
\begin{resolution}{Escalonamento}
\begin{columns}
\column{0.3\linewidth}
  \[\;\;\;
    \begin{amatrix}[2]{2}
      1 & 1 & 0  & 1 \\
      0 & 1 & -1 & 2
    \end{amatrix}
  \]
  \pause
\column{0.7\linewidth}
  Eliminamos os elementos acima do segundo pivô
  \\ \pause
  \(
    L_1 \leftarrow L_1 - L_2
  \)
\end{columns}
\vspace{-\baselineskip}
\[
  \pause
  L'_1
  \;=\;
  L_1 - L_2
  \pause
  \;=\;
  \begin{amatrix}[2]{2}
    1 & 1 & 0  & 1
  \end{amatrix}
  -
  \begin{amatrix}[2]{2}
    0 & 1 & -1 & 2
  \end{amatrix}
  \pause
  \;=\;
  \begin{amatrix}[2]{2}
    1 & 0 & 1  & -1
  \end{amatrix}
\]
\pause
\[
  \begin{amatrix}[2]{2}
    1 & 0 & 1  & -1 \\
    0 & 1 & -1 & 2
  \end{amatrix}
\]
\end{resolution}

%------------------------------------------------------------------------------%
\begin{resolution}{Solução}
  \[
    \begin{amatrix}[2]{2}
      1 & 0 & 1  & -1 \\
      0 & 1 & -1 & 2
    \end{amatrix}
  \]

  \pause
  Como a matriz tem posto completo

  \pause
  a inversa exite e pode ser lida do lado direito
  \pause
  \[
    A^{-1} =
    \begin{bmatrix}
      1  & -1 \\
      -1 & 2
    \end{bmatrix}
  \]
\end{resolution}

%------------------------------------------------------------------------------%
\begin{resolution}{Verificação}
  \onslide<+->{Podemos verificar o resultado calculando}
  \begin{align*}
    \onslide<+->{AA^{-1}}
    \onslide<+->{& =
      \begin{bmatrix}
        2 & 1 \\
        1 & 1
      \end{bmatrix}
      \begin{bmatrix}
        1  & -1 \\
        -1 & 2
      \end{bmatrix}}
    \\
    \onslide<+->{& =
      \begin{bmatrix}
        2\times 1 + 1 (-1) &  2(-1) + 1 \times 2\\
        1\times 1 + 1 (-1) &  1(-1) + 1 \times 2
      \end{bmatrix}}
    \\
    \onslide<+->{& =
      \begin{bmatrix}
        1 & 0 \\
        0 & 1
      \end{bmatrix}}
      \\
    \onslide<+->{& = I}
  \end{align*}
  \onslide<+->{Portanto, a matriz encontrada é realmente a inversa de $A$}
\end{resolution}

%------------------------------------------------------------------------------%
\endinput
